\documentclass{semantic-resume}

\begin{document}

\resumeheader
  {Dylan Madisetti}
  {Senior Software Engineer II, CoreWeave --- LLM agent systems, evaluation harnesses, and scientific computing}
  {San Francisco, CA \enspace|\enspace \href{mailto:contact@dylanmadisetti.com}{contact@dylanmadisetti.com} \enspace|\enspace \href{https://www.dylanmadisetti.com}{dylanmadisetti.com}}
  {\href{https://github.com/dmadisetti}{github.com/dmadisetti} \enspace|\enspace \href{https://linkedin.com/in/dylanmadisetti}{linkedin.com/in/dylanmadisetti}}

\section{Summary}
Senior software engineer and Mechanical Engineering PhD with 8+ years building
software across industry and research. Owns core architecture for marimo, a
23{,}000+ star scientific notebook project, spanning reusable Python workflows,
caching, and agent-facing correctness tooling. Background includes corresponding-author
simulation research, NASA C++/MPI systems, Google production NLU, and auditable LLM agent harnesses.

\section{Engineering Experience}
\resumeentry{CoreWeave}{2025--present}{Senior Software Engineer II --- marimo (acquired Oct. 2025)}
\begin{resumeitems}
  \item Authored \href{https://github.com/marimo-team/meps/blob/main/mep-0008.md}{MEP 0008}
    and implemented reusable-notebook architecture spanning top-level-function
    serialization, setup-cell protections, and notebook embedding.
  \item Designed and implemented marimo's caching stack, including persistent and
    in-memory decorators, cache introspection, custom serialization, cell-level
    cached execution, and signed cache bundles.
  \item Shipped \texttt{marimo check}, correctness tooling built for humans and coding
    agents that emits machine-readable diagnostics; drove it from public design
    through configuration and production rules.
  \item Authored the specification and security boundaries and built a proof of
    concept for \href{https://molab.marimo.io}{mo lab}'s sandboxed execution on
    CoreWeave Cloud; the product launched in public preview with GPU-backed sessions.
  \item Built a distributed-compute DAG prototype for marimo notebooks during a
    CoreWeave hackathon.
  \item Serve as a core maintainer of a 23{,}000+ star project (Nov. 2026): translate
    external user reports into shipped caching and pytest fixes, review contributors,
    publish releases, and mentor interns, junior colleagues, and SciPy participants.
\end{resumeitems}

\section{Agent Systems \& Evaluation}
\resumeentry{Cowboy}{2026--present}{Creator and maintainer}
\begin{resumeitems}
  \item Designed and open-sourced a Rust/Python/Nix runtime for persistent LLM agents
    with sandbox and controlled run-as boundaries, approval-gated effects,
    asynchronous message bridges, and audit-oriented execution records.
  \item Built and ran a prespecified four-task Harbor harness-validity calibration on
    Terminal-Bench 2.1 (1/4 task pass; 4/4 valid; zero infrastructure exclusions),
    validating a larger model-output cap and isolating a remaining runtime timeout.
\end{resumeitems}

\resumeentry{Google --- Mountain View}{Jul 2017--Jul 2018}{Software Engineer, Natural Language Understanding}
\begin{resumeitems}
  \item Developed and evaluated debiasing methods for word, sentence, and
    semantic-vector representations on Google's NLU research team.
  \item Developed components of a multi-tower semantic-model framework and trained
    TensorFlow graphs for deployment at scale.
  \item Built massively parallel C++/MapReduce analysis pipelines and
    hyperparameter-autotuning procedures for black-box optimization.
\end{resumeitems}

\resumeentry{NASA Langley Research Center}{Jun 2016--Aug 2016}{Research and Development Intern}
\begin{resumeitems}
  \item Developed custom C++ Distributed Point Source Method simulation software
    for massively parallel MPI clusters.
  \item Built a laser/vibrometer system for guided-wave nondestructive evaluation
    of damaged composites and automated experimental equipment with C\# serial I/O.
\end{resumeitems}

\section{Earlier Software Experience}
\small
Software Developer, Equiscript and BoomTown (2012--2013); Co-founder / Software
Developer, CloudCollege (2014--2015): shipped a beta-to-release product, ran user
research, and helped secure \$50{,}000 in seed funding.
\normalsize



\newpage
\section{Scientific Research}
\resumeentry{Johns Hopkins University}{2018--2026}{PhD Researcher; co-advised by Christopher D. Stiles, JHU APL}
\begin{resumeitems}
  \item Developed discrete-dislocation-dynamics and GPU-accelerated ray-tracing
    workflows to predict XRD peak broadening in additively manufactured 316L steel;
    corresponding author in \emph{Journal of the Mechanics and Physics of Solids}.
  \item Built \texttt{pd3}, an open-source C++/Python parallel dislocation-dynamics
    framework, and applied ML to reconstruct dislocation structures from XRD and 3D
    simulations; presented results at MRS and TMS.
  \item Developed hyperspectral characterization analysis with JHU APL for in-situ
    monitoring of laser powder-bed-fusion builds.
  \item NSF Graduate Research Fellow and teaching assistant for Machine Learning for
    Mechanical Engineers.
\end{resumeitems}

\resumeentry{Integrated Material Assessment \& Predictive Simulation Laboratory, USC}{2015--2017}{Undergraduate Researcher}
\begin{resumeitems}
  \item Developed CAD discretization and CUDA-accelerated point-cloud meshing for
    distributed-point-source and peridynamic simulations; also researched energy
    harvesting and acousto-elastic metamaterials, leading to a publication.
\end{resumeitems}

\section{Selected Publications and Talks}
\small
\begin{resumeitems}
  \item Madisetti et al. ``Computational investigation into XRD peak broadening
    effects with discrete dislocation dynamics.'' \emph{JMPS} 196 (2025),
    corresponding author.
  \item Lavenstein et al. ``The heterogeneity of persistent slip band nucleation
    and evolution in metals at the micrometer scale.'' \emph{Science} 370 (2020).
  \item Madisetti et al. ``Hash all the things: Caching for fast notebook restarts.''
    \emph{Proceedings of SciPy 2026}; presented ``Python Notebooks Are Better Now''
    at PyAI 2026.
\end{resumeitems}
\normalsize

\section{Technical Strengths}
\small
\resumeline{Languages \& systems}{Python, C++, Rust, C\#, SQL; Nix, WASM, MPI, CUDA, MapReduce, TensorFlow}
\resumeline{Scientific computing}{Discrete dislocation dynamics, GPU ray tracing, XRD, DPSM, peridynamics, CAD meshing}
\resumeline{Agent systems}{LLM harnesses, evaluations, sandboxing, approval gates, message bridges, distributed DAGs}
\normalsize

\section{Education}
\resumeline{Johns Hopkins University}{PhD, Mechanical Engineering, conferred 2026; NSF GRFP}
\resumeline{University of South Carolina}{BS, Mechanical Engineering, 2017; minors in Mathematics and Computer Science}

\end{document}
